% Intrinsic representation identities and exactness of the nonradial core.
\subsection{The eikonal core and its representation}
\label{subsec:nr-recovery}

Work on $A_{\mathrm{ess}}$, denoted by $A$, with nonradial weight operator.
By \cref{thm:nonradial-spectral} and \eqref{eq:nr-B-subalgebra},
\begin{equation}\label{eq:nr-recovery-splitting}
 A=B\oplus U,\qquad T|_B=\Id,\qquad T|_U=\tfrac7{27}\Id,
 \qquad B^2\subset B,\quad BU\subset U.
\end{equation}

\begin{proposition}[The eikonal core and its grading]\label{prop:nr-eikonal-core}
Assume $T|_{A_{\mathrm{ess}}}$ is not scalar.
Then \(H=0\), \(Z\ne0\), and the subalgebra \(B=E_T(1)\)
is eikonal. Moreover, $U^2\subset B$, so $A=B\oplus U$ is a
$\mathbb Z_2$-grading.
\end{proposition}

\begin{proof}
\Cref{lem:nr-negative-spaces}(i) and
\cref{thm:nonradial-spectral} give \(H=0\).
If \(Z=0\), the proof of \cref{thm:nonradial-spectral} would give
\(B=\R c\), so \(Z\ne0\). Fusion and invariance give
\[
 zz'=-\ip{z}{z'}c,\qquad ZW\subset W,\qquad
 w^2=\frac12|w|^2c+z_w,\quad z_w\in Z.
\]
For a unit \(z\in Z\), the self-adjoint map \(L_z\) preserves
\(W\). If \(w\in W\) is an eigenvector with eigenvalue \(y\),
then \(L_{e_\pm}w=(-1/4\pm\sqrt3y/2)w\), where
\(e_\pm=(-c\pm\sqrt3z)/2\). Since \(w\perp e_\pm\) and
\(Tw=w\), the quotient formula gives
\[
 p(-1/4+\sqrt3y/2)=p(-1/4-\sqrt3y/2)=1.
\]
The roots of \(p(t)-1\) are \(-1,-1/2,1/2\). Two such roots
with sum \(-1/2\) must be \(-1,1/2\), so \(y^2=3/4\).
Consequently \((L_z|_W)^2=3\Id/4\).
Scaling and the norm identity at \(w\in W\) yield
\[
 |zw|^2=\frac34|z|^2|w|^2,\qquad |z_w|^2=\frac34|w|^4.
\]
For \(b=tc+z+w\), its square has the orthogonal components
\[
 b^2=\left(t^2-|z|^2+\frac12|w|^2\right)c
          +(-2tz+z_w)+(tw+2zw).
\]
On taking norms, the terms linear in \(t\) cancel because
\(\ip{z}{z_w}=\ip{zw}{w}\). The displayed norm identities then
give \(|b^2|^2=|b|^4\). Differentiating within the metrized
subalgebra \(B\) proves \(b^3=|b|^2b\).

For the grading, use $U=P\oplus Q$ from \eqref{eq:nr-U-peirce}.
For the input eigenvalue pairs $(1/3,1/3)$, $(-1/6,-1/6)$,
and $(1/3,-1/6)$, the fusion rule \eqref{eq:nr-fusion}
permits, among the eigenvalues on $c^\perp$, respectively no output,
only $-1$, and only $1/2$.
Metric invariance therefore gives
\[
 p^2=\tfrac13|p|^2c,\qquad q^2\in\R c\oplus Z,\qquad pq\in W
 \qquad(p\in P,\ q\in Q).
\]
Polarization yields $U^2\subset B$; the other two grading relations
are already contained in \eqref{eq:nr-B-subalgebra}.
\end{proof}

At the maximal unit idempotent $c$, the decomposition is therefore
\[
 B=\R c\oplus Z\oplus W,\qquad U=P\oplus Q,
\]
with $L_c$-eigenvalues $1,-1,1/2,1/3,-1/6$, respectively, and $Z\ne0$.
To determine the action of $B$ on $U$, set $M_b=L_b|_U$ and
$\{A_1,A_2\}=A_1A_2+A_2A_1$.

\begin{proposition}\label{prop:nr-representation-identities}
For $b,d\in B$,
\begin{equation}\label{eq:nr-representation-quadratic}
\begin{aligned}
 M_{b^2}&=6M_b^2-\tfrac13|b|^2\Id_U,\\
 M_{bd}&=3\{M_b,M_d\}-\tfrac13\ip bd\Id_U.
\end{aligned}
\end{equation}
\begin{equation}\label{eq:nr-representation-cubic}
 108M_b^3-9|b|^2M_b-\ip{b^2}b\Id_U=0.
\end{equation}
\end{proposition}
\begin{proof}
Every unit idempotent $e\in B$ satisfies $\ip{e^2}e=1$, so it
attains the global spherical maximum in
\eqref{eq:nr-spherical-maximum}.
The maximality bounds and \cref{lem:nr-idempotent-quotient}, applied
on $U\subset e^\perp$, give
\[
 \operatorname{Spec}M_e\subset\{1/3,-1/6\}.
\]
Indeed, the third root $-7/6$ of $p(t)=7/27$ is below the lower
bound $-1$. Hence
$M_e=6M_e^2-\Id_U/3$.

Consider the homogeneous quadratic operator polynomial \(\mathscr Q(b)=M_{b^2}-6M_b^2+\tfrac13|b|^2\Id_U\).
Every nonzero $b\in B$ lies in an eikonal subalgebra of dimension
at most two. Indeed, differentiating $b^3=|b|^2b$ in direction $b^2$
gives
\[
 (b^2)^2=2\ip{b^2}b\,b-|b|^2b^2,
\]
so $\operatorname{span}\{b,b^2\}$ is such a subalgebra.
In dimension one it is generated by a unit idempotent, and hence
$\mathscr Q$ vanishes there.
In dimension two choose a positive spherical maximum $e$ for its cubic.
Since $|e^2|=1$, the multiplier equation gives $e^2=e$ without
rescaling. For a unit vector $z\perp e$, write
$ez=\alpha z$, $z^2=\alpha e+\beta z$.
Linearization at $e$ gives $2\alpha^2+\alpha=1$, while $z^3=z$
gives $\alpha\beta=0$ and $\alpha^2+\beta^2=1$.
Thus $\alpha=-1$, $\beta=0$, and
\[
 e,\qquad \frac{-e+\sqrt3z}{2},\qquad
 \frac{-e-\sqrt3z}{2}
\]
are three pairwise nonproportional unit idempotents.
A homogeneous quadratic polynomial on a plane vanishing on three
distinct lines is zero. Applying this to each matrix entry of
$\mathscr Q$ proves $\mathscr Q(b)=0$.
Polarization gives the second quadratic identity.
Finally, apply it to $b,d=b^2$. Since $b^3=|b|^2b$,
\[
 |b|^2M_b
 =3\{M_b,6M_b^2-\tfrac13|b|^2\Id_U\}
       -\tfrac13\ip{b^2}b\Id_U
 =36M_b^3-2|b|^2M_b-\tfrac13\ip{b^2}b\Id_U,
\]
which gives the cubic identity.
\end{proof}

The remaining identities on $U$ are as follows.

\begin{lemma}\label{lem:nr-U-cubic}
For $u\in U$, set $s=|u|^2$, $w=u^2\in B$, and $v=wu\in U$.
Then
\begin{equation}\label{eq:nr-U-identities}
 v=\tfrac s9u,\qquad |w|^2=\tfrac{s^2}9,\qquad
 w^2=\tfrac s3w,\qquad uv=\tfrac s9w.
\end{equation}
Moreover, for every $b\in B$,
\begin{equation}\label{eq:nr-mixed-polarization}
 -3wb+18u(bu)=sb.
\end{equation}
\end{lemma}
\begin{proof}
Pair the first identity in \eqref{eq:nr-representation-quadratic}
with $u$ and differentiate in $b$ to obtain
\eqref{eq:nr-mixed-polarization}.
At $u$, \eqref{eq:nr-vector-identity} reads
$4uv+w^2=7sw/9$.
Together with \eqref{eq:nr-mixed-polarization} at $b=w$, this gives
$uv=sw/9$ and $w^2=sw/3$.
Since $B$ is eikonal, $|w^2|^2=|w|^4$, and therefore
\[
 |u^2|^2\left(|u^2|^2-\tfrac19|u|^4\right)=0
 \qquad(u\in U).
\]
The first factor is not the zero polynomial: $U\ne0$ and the
only possible eigenvalues $1/3,-1/6$ of $M_c$ are nonzero, so
$\ip{u^2}c=\ip{M_cu}u$ is not identically zero.
The polynomial ring on $U$ is an integral domain, whence
$|u^2|^2=|u|^4/9$ identically.
Differentiating this identity within $U$ gives $4wu=4su/9$.
\end{proof}

\begin{proposition}[Exactness of the core]\label{prop:nr-core-exactness}
If the original cubic is irreducible over $\R$, then $B$ is exact,
\[
 \tr_B L_b=\tr_U M_b=0\qquad(b\in B),
\]
and $\dim B=\dim U+2$.
\end{proposition}
The proof is given in \cref{app:nr-core-exactness-proof}.
